% Cover letter (one page). Slots in [brackets] to be filled by the corresponding author.
\documentclass[11pt,a4paper]{article}
\usepackage[T1]{fontenc}
\usepackage[margin=22mm]{geometry}
\usepackage{parskip}
\pagestyle{empty}
\begin{document}
\noindent{}[Corresponding author name]\\{}
{}[Department, Institution, full postal address, Country]\\{}
[e-mail] \hfill [Date]

\medskip
\noindent To the Editor-in-Chief, \emph{Mechanical Systems and Signal Processing}

Dear Editor,

We submit the manuscript \textbf{``Tracking the support stiffness of railway wheelset bearings from
in-service axle-box acceleration without measuring rail roughness: a multi-pass debiased Whittle
likelihood''} for consideration as a research paper.

\textbf{Problem.} The vertical stiffness of the load path from a wheelset journal to its axlebox (roller
contacts, housing, adapter, pad) falls when fits loosen, when the housing or pad degrades, or when the
bearing clearance grows. Existing bearing monitoring detects temperature and impacts, not stiffness.
The support resonance lies in the band of axle-box acceleration that operators already record, but any
stiffness identified from it is confounded with the unmeasured rail roughness.

\textbf{Contributions.} To our knowledge, this is the first method to identify the vertical stiffness of
the axle-box bearing support from in-service axle-box acceleration without measuring rail roughness.
Passes at different speeds and payloads are pooled in one debiased Whittle likelihood in which the
roughness levels are free nuisance parameters. The forward model couples the instrumented wheelset to
its neighbours through the rail; ignoring that coupling biases the estimate by 5.6\%, far outside its
confidence interval.
For a representative freight wagon (Y25 bogies, ballasted UIC60 track) we give a complete simulation
error budget: Cram\'er--Rao bounds, Monte Carlo bias and coverage, model-form sensitivity, sweeps over
published parameter ranges, a check against nonlinear Hertzian contact, detection limits for changes
between monitoring periods, and false-alarm rates under track drift.

\textbf{Fit to MSSP.} The work combines system identification, statistical signal processing of
spectra and vehicle--track dynamics for condition monitoring, and builds on recent MSSP papers on
axle-box acceleration (Carrigan \& Talbot 2023; Shen et al.\ 2023). Suggested subject area:
[subject area as listed in the submission system].

The code that produces every number, table and figure is openly available
(\texttt{https://github.com/Shreyansh369/wheelset-s2pinn}; Zenodo DOI [to be added]). The manuscript has
not been published and is not under consideration elsewhere. All authors have approved the manuscript
and its submission to \emph{Mechanical Systems and Signal Processing}.

Yours sincerely,\\[2mm]
[Corresponding author name], on behalf of all authors
\end{document}
